\subsection{Strong Typing in Operations}

Python names are dynamically bound, but Python operations are strongly checked at runtime. This means that a name may be rebound to different kinds of objects, but an operation still requires objects that know how to participate in that operation.

Dynamic typing does not mean ``anything combines with anything''.

\begin{verbatim}
value = 42
value = "42"

print(f"value: {value!r}")
print(f"type:  {type(value)}")
\end{verbatim}

Possible output:

\begin{verbatim}
value: '42'
type:  <class 'str'>
\end{verbatim}

The name \texttt{value} was allowed to move from an integer object to a string object. But once it refers to a string object, string rules apply.

\subsubsection{Why Heterogeneous Operations Such as \texttt{"3" + 4} Fail Instead of Silently Converting}

The expression:

\begin{verbatim}
"3" + 4
\end{verbatim}

looks tempting, but Python rejects it.

A safe demonstration:

\begin{verbatim}
text_value = "3"
number_value = 4

try:
    result = text_value + number_value
    print(result)
except TypeError as err:
    print("operation failed")
    print(f"message: {err}")
    print(f"type(err): {type(err)}")
\end{verbatim}

Possible output:

\begin{verbatim}
operation failed
message: can only concatenate str (not "int") to str
type(err): <class 'TypeError'>
\end{verbatim}

The reason is that \texttt{+} does not have one universal meaning. For integers, \texttt{+} means numeric addition. For strings, \texttt{+} means concatenation.

\begin{verbatim}
print(3 + 4)
print("3" + "4")

print()
print(f"type(3 + 4):     {type(3 + 4)}")
print(f"type('3' + '4'): {type('3' + '4')}")
\end{verbatim}

Possible output:

\begin{verbatim}
7
34

type(3 + 4):     <class 'int'>
type('3' + '4'): <class 'str'>
\end{verbatim}

So the mixed expression is ambiguous:

\begin{verbatim}
"3" + 4
\end{verbatim}

Should Python produce the integer \texttt{7}? Should it produce the string \texttt{"34"}? Should it convert the integer to text, or the text to an integer?

Python refuses to guess.

This is strong typing in practice:

\begin{quote}
Python does not silently convert unrelated object types just to make an operation succeed.
\end{quote}

Inspect the two objects:

\begin{verbatim}
text_value = "3"
number_value = 4

print(f"text_value:   {text_value!r}")
print(f"number_value: {number_value!r}")

print()
print(f"type(text_value):   {type(text_value)}")
print(f"type(number_value): {type(number_value)}")

print()
print("selected str methods:")
print(f"{'upper':<12} {'upper' in dir(text_value)}")
print(f"{'replace':<12} {'replace' in dir(text_value)}")
print(f"{'__add__':<12} {'__add__' in dir(text_value)}")

print()
print("selected int methods:")
print(f"{'bit_length':<12} {'bit_length' in dir(number_value)}")
print(f"{'__add__':<12} {'__add__' in dir(number_value)}")
print(f"{'__mul__':<12} {'__mul__' in dir(number_value)}")
\end{verbatim}

Possible output:

\begin{verbatim}
text_value:   '3'
number_value: 4

type(text_value):   <class 'str'>
type(number_value): <class 'int'>

selected str methods:
upper        True
replace      True
__add__      True

selected int methods:
bit_length   True
__add__      True
__mul__      True
\end{verbatim}

Both objects have \texttt{\_\_add\_\_} behavior, but not for the same kind of partner. The string can concatenate with another string. The integer can add with another number. Python will not invent a hidden conversion between them.

More examples:

\begin{verbatim}
examples = [
    ('"3" + 4', '"3" + 4'),
    ('"D" * 3', '"D" * 3'),
    ('"D" - 1', '"D" - 1'),
]

for label, source_code in examples:
    print(f"trying: {label}")

    try:
        result = eval(source_code)
        print(f"result: {result!r}")
        print(f"type:   {type(result)}")
    except TypeError as err:
        print("TypeError")
        print(f"message: {err}")

    print()
\end{verbatim}

Possible output:

\begin{verbatim}
trying: "3" + 4
TypeError
message: can only concatenate str (not "int") to str

trying: "D" * 3
result: 'DDD'
type:   <class 'str'>

trying: "D" - 1
TypeError
message: unsupported operand type(s) for -: 'str' and 'int'
\end{verbatim}

The expression \texttt{"D" * 3} works because string repetition by an integer is a defined operation. The expression \texttt{"D" - 1} fails because subtracting an integer from a string has no defined meaning.

Strong typing does not mean that mixed-type operations are always forbidden. It means that mixed-type operations must be explicitly supported by the involved types.

\subsubsection{Explicit Conversion with \texttt{int()}, \texttt{float()}, \texttt{complex()}, \texttt{bool()}, and \texttt{str()}}

When the programmer wants conversion, the programmer must say so.

To add numerically:

\begin{verbatim}
text_value = "3"
number_value = 4

result = int(text_value) + number_value

print(f"result: {result}")
print(f"type:   {type(result)}")
\end{verbatim}

Possible output:

\begin{verbatim}
result: 7
type:   <class 'int'>
\end{verbatim}

To concatenate as text:

\begin{verbatim}
text_value = "3"
number_value = 4

result = text_value + str(number_value)

print(f"result: {result!r}")
print(f"type:   {type(result)}")
\end{verbatim}

Possible output:

\begin{verbatim}
result: '34'
type:   <class 'str'>
\end{verbatim}

The two programs give different results because they express different intentions.

Python provides built-in conversion functions:

\begin{verbatim}
int()
float()
complex()
bool()
str()
\end{verbatim}

A compact demonstration:

\begin{verbatim}
text = "42"

as_int = int(text)
as_float = float(text)
as_complex = complex(text)
as_bool = bool(text)
as_str = str(as_int)

print(f"original:   {text!r:<10} type={type(text)}")
print(f"int:        {as_int!r:<10} type={type(as_int)}")
print(f"float:      {as_float!r:<10} type={type(as_float)}")
print(f"complex:    {as_complex!r:<10} type={type(as_complex)}")
print(f"bool:       {as_bool!r:<10} type={type(as_bool)}")
print(f"str:        {as_str!r:<10} type={type(as_str)}")
\end{verbatim}

Possible output:

\begin{verbatim}
original:   '42'       type=<class 'str'>
int:        42         type=<class 'int'>
float:      42.0       type=<class 'float'>
complex:    (42+0j)    type=<class 'complex'>
bool:       True       type=<class 'bool'>
str:        '42'       type=<class 'str'>
\end{verbatim}

Conversions are not magic. They obey the rules of the target type.

This works:

\begin{verbatim}
print(int("42"))
print(float("3.14"))
print(complex("2+5j"))
\end{verbatim}

Possible output:

\begin{verbatim}
42
3.14
(2+5j)
\end{verbatim}

This fails:

\begin{verbatim}
try:
    number = int("3.14")
    print(number)
except ValueError as err:
    print("conversion failed")
    print(f"message: {err}")
    print(f"type(err): {type(err)}")
\end{verbatim}

Possible output:

\begin{verbatim}
conversion failed
message: invalid literal for int() with base 10: '3.14'
type(err): <class 'ValueError'>
\end{verbatim}

The string \texttt{"3.14"} is valid floating-point text, but it is not valid integer text.

A two-step conversion is possible if the programmer accepts the meaning:

\begin{verbatim}
text = "3.14"

as_float = float(text)
as_int = int(as_float)

print(f"as_float: {as_float}, type={type(as_float)}")
print(f"as_int:   {as_int}, type={type(as_int)}")
\end{verbatim}

Possible output:

\begin{verbatim}
as_float: 3.14, type=<class 'float'>
as_int:   3, type=<class 'int'>
\end{verbatim}

This does not round. It truncates toward zero.

\begin{verbatim}
values = [3.99, -3.99, 42.0]

for value in values:
    converted = int(value)
    print(f"int({value}) = {converted}")
\end{verbatim}

Possible output:

\begin{verbatim}
int(3.99) = 3
int(-3.99) = -3
int(42.0) = 42
\end{verbatim}

The function \texttt{bool()} converts by truth-value testing:

\begin{verbatim}
values = [0, 1, 42, "", "0", "False", None]

for value in values:
    converted = bool(value)
    print(f"{repr(value):<8} -> {converted:<5} type={type(converted)}")
\end{verbatim}

Possible output:

\begin{verbatim}
0        -> False type=<class 'bool'>
1        -> True  type=<class 'bool'>
42       -> True  type=<class 'bool'>
''       -> False type=<class 'bool'>
'0'      -> True  type=<class 'bool'>
'False'  -> True  type=<class 'bool'>
None     -> False type=<class 'bool'>
\end{verbatim}

The strings \texttt{"0"} and \texttt{"False"} are true because they are non-empty strings. \texttt{bool()} does not parse English meaning. It checks truthiness.

The function \texttt{str()} creates a text representation:

\begin{verbatim}
values = [42, 3.14, True, None, 2 + 5j]

for value in values:
    text = str(value)
    print(f"{repr(value):<10} -> {text!r:<10} type={type(text)}")
\end{verbatim}

Possible output:

\begin{verbatim}
42         -> '42'       type=<class 'str'>
3.14       -> '3.14'     type=<class 'str'>
True       -> 'True'     type=<class 'str'>
None       -> 'None'     type=<class 'str'>
(2+5j)     -> '(2+5j)'   type=<class 'str'>
\end{verbatim}

The practical rule is:

\begin{quote}
If the program needs a conversion, write the conversion explicitly. Do not depend on Python to guess whether text should become a number or a number should become text.
\end{quote}

\subsubsection{Numeric Promotion Boundaries: Integer, Float, Complex, and Boolean Interactions}

Python does allow many mixed numeric operations. These are not silent conversions between unrelated domains. They are defined interactions inside Python's numeric tower.

A simple table:

\begin{verbatim}
values = [
    ("int + int",       40 + 2),
    ("int + float",     40 + 2.5),
    ("int + complex",   40 + 2j),
    ("float + complex", 3.5 + 2j),
    ("bool + int",      True + 41),
    ("bool + float",    True + 0.5),
]

for label, result in values:
    print(f"{label:<18} -> {result!r:<10} type={type(result)}")
\end{verbatim}

Possible output:

\begin{verbatim}
int + int          -> 42         type=<class 'int'>
int + float        -> 42.5       type=<class 'float'>
int + complex      -> (40+2j)    type=<class 'complex'>
float + complex    -> (3.5+2j)   type=<class 'complex'>
bool + int         -> 42         type=<class 'int'>
bool + float       -> 1.5        type=<class 'float'>
\end{verbatim}

The broad pattern is:

\begin{verbatim}
int + int       -> int
int + float     -> float
int + complex   -> complex
float + complex -> complex
bool + int      -> int
bool + float    -> float
\end{verbatim}

This happens because Python defines these numeric combinations. The result type is chosen so that it can represent the result of the operation.

Inspect the operands:

\begin{verbatim}
a = 42
b = 0.5
c = 2j
d = True

print(f"a: {a!r:<8} type={type(a)}")
print(f"b: {b!r:<8} type={type(b)}")
print(f"c: {c!r:<8} type={type(c)}")
print(f"d: {d!r:<8} type={type(d)}")
\end{verbatim}

Possible output:

\begin{verbatim}
a: 42       type=<class 'int'>
b: 0.5      type=<class 'float'>
c: 2j       type=<class 'complex'>
d: True     type=<class 'bool'>
\end{verbatim}

Now combine them:

\begin{verbatim}
print(f"a + b: {a + b!r:<8} type={type(a + b)}")
print(f"a + c: {a + c!r:<8} type={type(a + c)}")
print(f"b + c: {b + c!r:<8} type={type(b + c)}")
print(f"d + a: {d + a!r:<8} type={type(d + a)}")
\end{verbatim}

Possible output:

\begin{verbatim}
a + b: 42.5     type=<class 'float'>
a + c: (42+2j)  type=<class 'complex'>
b + c: (0.5+2j) type=<class 'complex'>
d + a: 43       type=<class 'int'>
\end{verbatim}

This is still strongly typed. Python is not guessing that arbitrary objects should convert. It is applying defined numeric rules.

For example, \texttt{int} and \texttt{str} are not part of one shared numeric operation:

\begin{verbatim}
try:
    result = 42 + "1"
except TypeError as err:
    print("operation failed")
    print(f"message: {err}")
\end{verbatim}

Possible output:

\begin{verbatim}
operation failed
message: unsupported operand type(s) for +: 'int' and 'str'
\end{verbatim}

Similarly, \texttt{complex} supports arithmetic but not ordering:

\begin{verbatim}
try:
    print((3 + 4j) < 10)
except TypeError as err:
    print("comparison failed")
    print(f"message: {err}")
\end{verbatim}

Possible output:

\begin{verbatim}
comparison failed
message: '<' not supported between instances of 'complex' and 'int'
\end{verbatim}

Equality is allowed:

\begin{verbatim}
print(f"3 + 0j == 3:   {3 + 0j == 3}")
print(f"3.0 == 3:      {3.0 == 3}")
print(f"True == 1:     {True == 1}")
print(f"False == 0:    {False == 0}")
\end{verbatim}

Possible output:

\begin{verbatim}
3 + 0j == 3:   True
3.0 == 3:      True
True == 1:     True
False == 0:    True
\end{verbatim}

This is numeric equality, not type identity.

\begin{verbatim}
print(f"type(3 + 0j): {type(3 + 0j)}")
print(f"type(3):      {type(3)}")

print()
print(f"(3 + 0j) == 3:      {(3 + 0j) == 3}")
print(f"type(3 + 0j) is int: {type(3 + 0j) is int}")
\end{verbatim}

Possible output:

\begin{verbatim}
type(3 + 0j): <class 'complex'>
type(3):      <class 'int'>

(3 + 0j) == 3:      True
type(3 + 0j) is int: False
\end{verbatim}

Booleans require special care because \texttt{bool} is a subclass of \texttt{int}.

\begin{verbatim}
flag = True

print(f"flag: {flag}")
print(f"type(flag): {type(flag)}")

print()
print(f"isinstance(flag, bool): {isinstance(flag, bool)}")
print(f"isinstance(flag, int):  {isinstance(flag, int)}")

print()
print(f"flag + 41: {flag + 41}")
print(f"type(flag + 41): {type(flag + 41)}")
\end{verbatim}

Possible output:

\begin{verbatim}
flag: True
type(flag): <class 'bool'>

isinstance(flag, bool): True
isinstance(flag, int):  True

flag + 41: 42
type(flag + 41): <class 'int'>
\end{verbatim}

This works, but it is often bad style unless the programmer is deliberately counting truth values.

Acceptable:

\begin{verbatim}
passed_test_1 = True
passed_test_2 = False
passed_test_3 = True

passed_count = passed_test_1 + passed_test_2 + passed_test_3

print(f"passed_count: {passed_count}")
\end{verbatim}

Possible output:

\begin{verbatim}
passed_count: 2
\end{verbatim}

Less clear:

\begin{verbatim}
total = True + 41
print(total)
\end{verbatim}

Possible output:

\begin{verbatim}
42
\end{verbatim}

The practical numeric promotion rule is:

\begin{quote}
Python defines mixed operations among its numeric types, commonly moving from \texttt{bool}/\texttt{int} toward \texttt{float} and then toward \texttt{complex} when needed. This is not the same as arbitrary silent conversion between unrelated types such as \texttt{str} and \texttt{int}.
\end{quote}

The operational summary is:

\begin{verbatim}
"3" + 4          -> TypeError
int("3") + 4     -> 7
"3" + str(4)     -> "34"

42 + 0.5         -> 42.5
42 + 2j          -> (42+2j)
True + 41        -> 42
\end{verbatim}

Names are flexible. Operations are checked. Objects keep their runtime type rules.